\section{Related Work}\label{sec:related}

The most useful way to locate this paper is to start from the \emph{specific gap} that existing activation-steering methods share. SEKA~\citep{SEKA} and AdaSEKA~\citep{AdaSEKA} learn a low-rank projector from contrastive QA pairs and inject it into attention keys, achieving state-of-the-art single-concept editing. ITI~\citep{ITI} and CAA~\citep{CAA} add activation directions to the residual stream to shift behavior distributions. PASTA~\citep{PASTA} directly rescales post-softmax attention weights, and Focus Directions~\citep{FocusDirections} jointly edits keys and queries to sharpen long-context attention. The shared appeal is that all of them intervene at inference time without further fine-tuning. But a second feature is also shared: \emph{none of these methods carries an explicit state about which tools have already been emitted}. Stationary key-side amplification, residual-stream directions, post-softmax reweighting, and focus directions are all fixed transformations independent of the decoding step. Theorem~\ref{thm:no-memory} pins down why this common property becomes a limit on multi-tool tasks, and the paper uses this fact to surface the design space those methods do not occupy --- the sign-flipped query contraction.

Our separation from the prior line reduces to a single underlying choice: ``do not amplify keys, contract queries instead.'' The task differs --- multi-tool selection requires two or more distinct tools in sequence, where ``re-amplifying an already chosen direction'' is actively harmful. The operator location and sign differ --- we subtract from queries with a negative coefficient, and a negative Q-side projection acts as a coarse coverage prior (Theorem~\ref{thm:q-sign}). The basis source differs --- we derive $B_{\mathrm{ont}}$ from the tool catalog's four facet labels (domain, function action, io type, tool category) via Gram-Schmidt rather than from contrastive QA pairs.

Benchmark choices are narrow and deliberate. MetaTool Subtask4~\citep{MetaTool} assigns exactly two ground-truth tools to each of 497 prompts, the cleanest benchmark for sequential coverage. $\tau^2$-bench retail~\citep{tau2bench} has tool counts from 0 to 13, so short and long sequence regimes coexist in one benchmark, enabling a measurable regime-split. BFCL~\citep{BFCL} parallel/multiple subsets provide cross-validation. KV-cache compression~\citep{KIVI} is orthogonal to our focus. All experiments use the Qwen2.5-Instruct family (7B for the main body; 1.5B/3B/14B only in the appendix size sweep)~\citep{qwen25}.
